    \def\stat{zat-kis}



\def\tit{НЕКОТОРЫЕ ПОДХОДЫ К ФОРМИРОВАНИЮ

    НОРМАТИВНО-ТЕХНИЧЕСКОЙ БАЗЫ ДЛЯ~СОЗДАНИЯ ЕДИНОГО

ИНФОРМАЦИОННОГО ПРОСТРАНСТВА РОССИИ}



\def\titkol{Некоторые подходы к формированию

    нормативно-технической базы для~создания ЕИП} %единого информационного пространства   России}



\def\autkol{А.\,А. Зацаринный, Э.\,В. Киселев}



\def\aut{А.\,А. Зацаринный$^1$, Э.\,В.~Киселев$^2$}



\titel{\tit}{\aut}{\autkol}{\titkol}



%{\renewcommand{\thefootnote}{\fnsymbol{footnote}}

%\footnotetext[1]

%{Работа выполнена при частичной финансовой поддержке  программы

%<<Интеллектуальные информационные технологии, системный анализ и

%автоматизация>> (проект 1.7).}}







\renewcommand{\thefootnote}{\arabic{footnote}}

\footnotetext[1]{Институт проблем информатики Российской академии наук,

     azatsarinny@ipiran.ru}

\footnotetext[2]{Институт проблем информатики Российской академии наук,

     ekiselev@ipiran.ru}





\label{st\stat}



                  \thispagestyle{headings}







\Abst{На основе существующих нор\-ма\-тив\-но-пра\-во\-вых документов

федерального уровня сформулированы цели формирования и развития единого

информационного пространства (ЕИП) Российской Федерации

(ЕИП РФ). Определен обобщенный состав

ЕИП на основе совокупности информационных ресурсов (ИР)

и~информационной инфраструктуры (ИИ) страны. Введено понятие участников этого

пространства как основных источников и~потребителей ИР.

Выявлены три основные проблемы создания и~использования

ЕИП с~учетом реального состояния информатизации страны. На основе системного

подхода определены основные принципиальные положения по созданию и~использованию

ЕИП. Сформулирован подход и~даны предложения по

документированию ЕИП как мегасистемы, для которой

предложено семь качественно новых типов документов. Результаты статьи являются

постановочными для дальнейшего рассмотрения вопросов, связанных с формированием

и~взаимодействием ИР в~со\-ста\-ве ЕИП~РФ.}



\KW{единое информационное пространство Российской Федерации; информационные

ресурсы; информационная инфраструктура; централизация; интеграция; информационная

безопасность; телекоммуникационное обеспечение; системный подход; мегасистема;

нор\-ма\-тив\-но-тех\-ни\-че\-ские документы}



\DOI{10.14357\!/\!08696527140413}







\vskip 12pt plus 9pt minus 6pt





      \thispagestyle{headings}





     \section{Вводные положения}



     Весь мир, и в том числе Россия, находится в состоянии перехода от

<<индустриального общества>> к~<<обществу информационному>>.

Информация и~ИР становятся экономическими

категориями, одним из основных национальных ресурсов.



     В нашей стране актуальность, необходимость, цели и основные задачи

создания ЕИП РФ определены целым рядом нор\-ма\-тив\-но-пра\-во\-вых документов на

федеральном уровне~[1--7]. Так, в качестве основных предусмотрены

следующие цели формирования и развития ЕИП РФ~[1--4, 6]:

\begin{itemize}

\item  обеспечение прав граждан на информацию, провозглашенных

Конституцией Российской Федерации;

     \item создание и поддержание уровня информационного потенциала,

необходимого для устойчивого развития российского общества;

     \item обеспечение взаимоувязанности и согласованности решений,

принимаемых федеральными органами государственной власти, органами

власти субъектов РФ и~органами местного самоуправления (ОГВ);

     \item повышение уровня правосознания граждан путем предо\-став\-ления

им свободного доступа к~правовым и~нормативным документам,

определяющим их права, обязанности и возможности;

     \item обеспечение контроля со стороны граждан и~общественных

организаций за деятельностью ОГВ;

     \item повышение деловой и общественной активности граждан путем

предо\-став\-ления равной с~государственными структурами возможности

пользоваться открытой на\-уч\-но-тех\-ни\-че\-ской,

     со\-ци\-аль\-но-эко\-но\-ми\-че\-ской, об\-щест\-вен\-но-по\-ли\-ти\-че\-ской

     информацией, а~также информационными фондами сферы образования,

культуры и~т.\,д.;

     \item интеграцию с европейским и мировым информационным

пространством.

     \end{itemize}



     Для достижения поставленных целей в~\cite{1-zk} ЕИП РФ

представляется как совокупность баз и банков данных, технологий их ведения и

использования, ин\-фор\-ма\-ци\-он\-но-те\-ле\-ком\-му\-ни\-ка\-ци\-он\-ных

сис\-тем (ИТКС) и~сетей, функционирующих на основе единых принципов

совместимости и~по общим правилам, обеспечивающим информационное

взаимодействие организаций и~граждан, а~также удовлетворение их

информационных потребностей.



     Обобщенный состав ЕИП РФ складывается из двух укрупненных

компонентов~\cite{1-zk}:

     \begin{enumerate}[(1)]

     \item ИР, содержащих данные, сведения

и~знания, зафиксированные на соответствующих носителях информации.

В~ИР входят информация, информационное общество, информационный

потенциал общества, информационная культура;

     \item ИИ, включающей

организационные структуры, которые обеспечивают функционирование

и~развитие ЕИП РФ, а~так\-же ап\-па\-рат\-но-про\-грам\-мные средства,

обеспечивающие доступ организаций и~граждан к~ИР

на основе соответствующих информационных технологий (ИТ).

     \end{enumerate}



     Различные организации, структуры и~т.\,п.\ (далее обобщенно~---

организации), имеющие собственные ИР и~осуществляющие информационное

взаимодействие посредством ИИ с~другими организациями и~физическими

лицами в~пределах ЕИП, ниже

определяются как Участники ЕИП РФ. Многочисленные Участники являются

основными источниками и~потребителями ИР

     в~со\-ста\-ве ЕИП РФ. Между Участниками в~рамках ЕИП РФ

осуществляется информационное взаимодействие по правилам ЕИП РФ.



     \section{Основные проблемы создания единого информационного пространства России}



     Формирование ЕИП РФ и соответствующих государственных

ИР представляет собой \textit{глобальную задачу}~---

общегосударственную, межотраслевую и~межрегиональную. Она требует

решения сложных организационных и~тех\-ни\-ко-тех\-но\-ло\-ги\-че\-ских

вопросов, значительных затрат и~не может быть решена одномоментно. При

этом необходим комплексный учет со\-ци\-аль\-но-эко\-но\-ми\-че\-ских,

правовых и~политических аспектов информатизации общества, всестороннее

использование организационного, технологического, технического

и~нормотворческого опыта развития информационных пространств ведущих

стран.



     Из всего множества проблем следует выделить \textit{три ключевых}.



     \textit{Первая проблема}~--- совершенствование нормативной базы,

обеспечение единства нор\-ма\-тив\-но-пра\-во\-вой,

     нор\-ма\-тив\-но-тех\-ни\-че\-ской и~нор\-ма\-тив\-но-спра\-воч\-ной

информации и~централизованное, взаимоувязанное управление этими

составляющими в~рамках ЕИП РФ. В~том числе необходимо разработать

     эко\-но\-ми\-ко-пра\-во\-вые основы:

     \begin{itemize}

     \item законодательные и нормативные акты, определяющие права

и~обязанности юридических и~физических лиц по формированию

и~использованию ИР, средств их обработки

и~доставки;

     \item экономические регуляторы, обеспечивающие стимулирование

активного формирования и~использования ИР.

     \end{itemize}



     \textit{Вторая проблема}~--- централизованное формирование ИР

и~обеспечение оперативного доступа к~ним в~рамках ЕИП РФ. Необходимо

провести весьма существенные работы по включению ИР в~ЕИП РФ,

предоставление их на законном основании органам управления

государственной власти, хозяйствующим субъектам и гражданам.



     \textit{Третья проблема}~--- обеспечить соблюдение принципов

системного подхода к~информатизации при построении ЕИП РФ. Как

показывает анализ, такой подход практически игнорируется многими

потенциальными Участниками вследствие сложившейся практики так

называемой <<лоскутной>> информатизации, когда у~Участника имеется

множество самостоятельных, зачастую уникальных информационных систем,

не совместимых между собой, не приспособленных к~совершенствованию

и~развитию, дорогостоящих, трудных в~эксплуатации.



     \section{Основные системотехнические положения по созданию

и~использованию единого информационного пространства России}



     На основе углубленного системного анализа существующей

     нор\-ма\-тив\-но-пра\-во\-вой базы, других документов и

     материалов~[1--21] сформулированы следующие системотехнические

положения по созданию и использованию ЕИП РФ.

     \begin{enumerate}[1.]

     \item Единое информационное пространство России должно

     создаваться как конгломерация

общенационального масштаба, соединяющая в~одно целое \mbox{отдельные}

разнообразные автоматизированные системы (АС, автоматизированные информационные

сис\-те\-мы (АИС), ИТКС, информационно-управ\-ля\-ющие сис\-те\-мы и~др.)\

России на основе паритетности~--- при сохранении этими системами своих

черт, свойств и~ИР. К~тому же ЕИП РФ должно входить на основе

паритетности в~мировое информационное пространство. Таким образом, ЕИП

РФ с полным основанием можно считать <<Сис\-те\-мой сис\-тем>>

и~подобную систему рассматривать как мегаинфраструктурное образование~---

\textit{мегасистему}.



     Данной мегасистеме свойственны следующие основные характерные,

специфические особенности: глобальность, множественность Участников,

постоянные изменения состава ЕИП РФ, перманентное видоизменение самих

ИР, многочисленные изменения нор\-ма\-тив\-но-пра\-во\-вой базы ЕИП РФ.

     \item Единое информационное пространство России

     следует создавать на основе двух основополагающих принципов:

     \begin{itemize}

\item[(а)] \textit{адекватность} перспективного ЕИП РФ потребностям

функциональной деятельности Участников~--- всех организаций России,

входящих в~это пространство, с~учетом расширения ЕИП РФ. Это означает,

что ЕИП РФ целесообразно создавать с~некоторым опережением

потенциальных функциональных потребностей Участников ЕИП РФ;

\item[(б)] \textit{органичное включение Участников~--- конкретных

организаций РФ~--- в ЕИП РФ}.

В~значительной мере этому должна способствовать уже действующая

государственная программа Российской Федерации <<Информационное

общество (2011--2020~годы)>>~[6], определяющая конкретные организации

РФ (в~части их касающейся), которые должны войти на правах Участников

в~создаваемое <<Информационное общество>> России. Соответственно,

Участники должны быть подключены к~реализации этой программы, в~том

числе и~те, которые охвачены программой электронного правительства

РФ~--- в~час\-ти, касающейся~[5, 10, 11, 13--17]. Каждый Участник формирует

ИР именно для ЕИП РФ, возможно, эти ИР являются

частью его внутренних ресурсов.

\end{itemize}

     \item  Единое информационное пространство России следует создавать на

основе органичного сочетания \textit{централизации и интеграции}, прежде

всего, при территориальной распределенности Участников (например, по

регионам страны, возможно зарубежных объектов этих организаций).



     \textit{Централизация} в составе ЕИП РФ предусматривает перенесение

на высший уровень Участника (на федеральном или региональном уровне)

выполнения всех основных информационных процессов (сбора, накопления,

обработки, хранения, архивирования, защиты информации и~др.).



     Естественно, централизация в таком виде касается только той части

ИР Участника, которые подготавливаются им для

использования другими Участниками в~составе ЕИП РФ.



     \textit{Интеграция} в составе ЕИП РФ ориентирует Участника на

органичное объединение у~него существующих и вновь разрабатываемых

АИС в единую систему~--- ЕИС,

с~созданием ЕИП в~пределах Участника,

которое должно быть органично совместимым с~ЕИП РФ. При этом

осуществляется рациональное структурирование компонентов ЕИС как единой

сис\-те\-мы-ин\-те\-гра\-то\-ра и~взаимодействие компонентов между собой

и~с~внешними системами в рамках ЕИП РФ. Рациональное структурирование

при интегрировании выполняется с~учетом наличия существующей,

унаследованной ИИ. Функциональность

унаследованных АИС в~со\-ста\-ве единой системы Участника сохраняется, но

может видоизменяться ее реализация. Это означает, что вместо использования

самостоятельных существующих АИС их функциональность выполняется

в~со\-ста\-ве единой сис\-те\-мы-ин\-тегра\-то\-ра как задачи и (или) комплексы

функциональных задач на высшем уровне данной организации~--- Участника

ЕИП РФ.



     Результаты органичного сочетания централизации и интеграции в единой

сис\-те\-ме-интеграторе ориентируются на обеспечение:

     \begin{itemize}

\item достаточной гибкости структуры ЕИС в~процессе ее поэтапного создания,

развития и~модернизации;

\item относительной простоты про\-грам\-мно-ап\-па\-рат\-ных комплексов,

используемых на объектах массового применения;

\item работы ЕИС в режиме, близком к~режиму реального времени;

\item высокой надежности функционирования ЕИС в~целом;

\item максимально возможного сокращения совокупной стоимости владения

единой системой;

\item сравнительно быстрого, полного и~относительно безболезненного

перехода от имеющего место существующей <<лоскутности>>

информатизации в~структурах РФ при множественности АИС к~рационально

отструктурированной единой сис\-те\-ме-инте\-гра\-тору;

\item эффективного доступа Участника в ЕИП РФ посредством своей

ЕИС-ин\-тегратора.

\end{itemize}

     \item  В ЕИП РФ должен постоянно осуществляться автоматизированный

функциональный мониторинг и~выполнение соответствующих

     учет\-но-ана\-ли\-ти\-че\-ских работ, необходимых для обеспечения

полноценной функциональной деятельности Участников~--- организаций РФ

в~рамках ЕИП.



     Это означает, что соответствующие единые

     сис\-те\-мы-ин\-те\-гра\-то\-ры должны участвовать

в~автоматизированном функциональном мониторинге в~со\-ста\-ве ЕИП

     РФ~--- в~час\-ти их касающейся. При этом с~использованием ЕИП РФ

сфера выполнения аналитических и~учетных работ может быть существенно

расширена.

     \item В составе ЕИП РФ должен быть обеспечен эффективный

информационный обмен~--- на качественно новом уровне, с~использованием

новейших технологий.



     Главную роль при реализации информационного обмена должны

выполнять высокопроизводительные, территориально распределенные

телекоммуникационные структуры (сети, системы), которые доходят до

каждого объекта всех Участников ЕИП РФ.

     \item В составе ЕИП РФ должна быть обеспечена юридическая

значимость электронной информации.



     Достижение юридической значимости электронной информации

включает комплекс мер и решений современного уровня, рассчитанных на

обеспечение необходимой степени защищенности разнообразных

ИР и~взаимодействия Участников в~составе ЕИП РФ.

     \item В ЕИП РФ необходимо обеспечить непрерывный контроль

состояния и~управ\-ле\-ние функционированием системы в~целом и~ее

компонентов, включая централизованное управление конфигурацией. При этом

охватываются эксплуатационные процессы, оптимизационные процессы,

процессы развертывания и~управления изменениями, процессы поддержки

и~управления особыми ситуациями.

     \item В ЕИП РФ должно осуществляться комплексное обеспечение

информационной безопасности. Особое внимание следует постоянно уделять

комплексной защите государственных и~негосударственных

ИР, входящих в~ЕИП РФ.



     Формирование ЕИП РФ и взаимодействие в~со\-ста\-ве ЕИП РФ

осуществляется на основе обязательного принципа, а~именно: все без

исключения функционально ориентированные ИР

защищены надежно и~должным образом. Это означает, что \textit{ЕИП РФ

оперирует исключительно защищенными ИР (ЗИР)}.



     Обеспечение информационной безопасности в~со\-ста\-ве ЕИП РФ

преду\-смат\-ри\-вает использование комплекса взаимоувязанных мер и~решений

современного уровня. В~том числе могут использоваться для необходимых

функ\-ци\-о\-наль\-но-ори\-ен\-ти\-ро\-ван\-ных ЗИР различные контуры:

открытый (О-кон\-тур); конфиденциальный (К-кон\-тур), секретный

     (С-кон\-тур).

     \item В ЕИП РФ должно осуществляться непрерывное комплексное

обеспечение эксплуатации, модернизации и~сопровождения данной

мегасистемы в целом и ее компонентов в течение жизненного цикла.



     Комплексное обеспечение осуществляется на основе использования

профильных государственных стандартов и~лучших практик, в~том числе

с~использованием соответствующих разработанных и~утвержденных

регламентов, и~непрерывное отслеживание процессов эксплуатации,

модернизации и~сопровождения на соответствующих стадиях. Необходимо

ввести обще\-сис\-тем\-ный компонент для централизованного управления

и~контроля процессов эксплуатации, модернизации и сопровождения в~составе

ЕИП~РФ.

     \item На всех стадиях и этапах создания ЕИП РФ и его компонентов

(составных частей) должно осуществляться документирование данной

мегасистемы. Особое внимание следует уделять тщательной отработке

эксплуатационной документации, прежде всего достаточности состава

и~содержания документов, изложения материалов и~качества документов,

удобства использования их при эксплуатации.



     Известно, что детальное документирование многих АС, которые

потенциально должны войти в~ЕИП РФ, при их создании и~эксплуатации

осуществляется в~соответствии с~требованиями и методическими указаниями

стандартов, например комплекса стандартов на автоматизированные системы

ГОСТ серии~34. Такой подход, к~сожалению, не применим непосредственно

к~такому специфическому инфраструктурному образованию, как ЕИП РФ.

Учитывая важность документирования, ниже излагается подход

и~предложения по этой части.

     \item Особое внимание в ЕИП РФ должно быть уделено обеспечению

единства нормативной информации трех видов~--- правовой (НПИ),

технической (НТИ), справочной (НСИ)~--- и централизованное управление ими

в~рамках ЕИП.



     Это означает, что в составе ЕИП РФ должна быть создана и~постоянно

использоваться централизованная система нормативной информации

как один из важнейших общесистемных компонентов данной мегасистемы.

Централизованная система нормативной информации должна иметь

в~своем составе соответственно подсистемы трех видов~--- ПНПИ, ННТИ, ПНСИ.

     \end{enumerate}



     Завершая рассмотрение принципиальных положений по созданию

и~использованию ЕИП РФ, следует обратить внимание на особую значимость

решения организационных вопросов. Речь идет об определении

соответствующих органов государственной власти с~правами создания ЕИП

РФ, координации формирования и развития всех видов государственных

ИР и~включения в~ЕИП негосударственных ИР. Ответственными

уполномоченными органами, например, можно было бы определить:

     \begin{itemize}

\item Администрацию Президента РФ в~час\-ти функциональных задач,

выпол\-ня\-емых ЕИП РФ;

\item Федеральную службу охраны России в части обеспечения функционирования и~эксплуатации

инф\-ра\-струк\-ту\-ры ЕИП РФ на всех стадиях жизненного цикла;

\item Минкомсвязи России в части телекоммуникационного обеспечения ЕИП

РФ;

\item Федеральную службу безопасности

и Федеральную службу по техническому и~экспортному

контролю России в~час\-ти комплексного обеспечения

информационной безопасности в~со\-ста\-ве ЕИП РФ.

\end{itemize}



     \section{Подход и~предложения по~документированию на~уровне

единого информационного пространства России}



     Мегасистему ЕИП РФ, как и любую систему, необходимо

целенаправленно проектировать, создавать и~правильно эксплуатировать.



     В настоящем разделе с учетом специфики ЕИП РФ формулируется

подход и~даются предложения по документированию данной мегасистемы.

Предлагаются семь качественно новых типов документов более высокого,

обобщающего уровня:

     \begin{enumerate}[1.]

\item Стратегия создания и использования инфраструктуры ЕИП РФ.

\item Тезаурус инфраструктуры ЕИП РФ.

\item Совокупность стандартов в части инфраструктуры ЕИП РФ.

\item Описание процессов Участника в~со\-ста\-ве ЕИП РФ,

соглашение с Участником об уровне предоставления услуг

и~регламенты.

\item Схема информатизации Участника ЕИП РФ.

\item Схема информационных потоков.

\item Схема взаимной зависимости ИР

Участников ЕИП РФ и~их сервисов.

\end{enumerate}



     Краткое содержание документов дано в таблице.



     \begin{table}[b]\small

     \begin{center}

      \begin{tabular}{|p{26mm}|p{95mm}|}

     \multicolumn{2}{c}{Типы документов с их кратким содержанием}\\[6pt]

      \hline

\multicolumn{1}{|c|}{Тип документа}&\multicolumn{1}{c|}{Краткое содержание документа}\\

\hline

1. Стратегия создания и~использования инфраструктуры ЕИП РФ&

Цели и задачи

создания и использования инфраструктуры ЕИП РФ. Подход к организационному

построению инфраструктуры ЕИП РФ. Общая интеграционная архитектура

функционирования ЕИП РФ. Подход, система приоритетов и~правил построения

инфраструктуры ЕИП РФ. Ключевые решения со стороны Участников в~части

информатизации с~учетом вхож\-де\-ния их в~ЕИП РФ на принципах интеграции. Принципы

инфраструктурного взаимодействия Участников в рамках ЕИП РФ с использованием

ЗИР при взаимодействии. Определение стратегически

важных сбалансированных показателей эффективности инфраструктуры ЕИП РФ. Основные

планы развития инфраструктуры ЕИП РФ в~краткосрочной и среднесрочной перспективе,

бюджетная и кадровая политика реализации планов\\

\hline

\multicolumn{2}{r}{\textit{Окончание таблицы на с.~\pageref{ppa}}}

\end{tabular}

\end{center}

\end{table}



\begin{table}\small

     \begin{center}

     \label{ppa}

      \begin{tabular}{|p{26mm}|p{95mm}|}

     \multicolumn{2}{c}{Типы документов с их кратким содержанием (\textit{окончание})}\\[6pt]

      \hline

\multicolumn{1}{|c|}{Тип документа}&\multicolumn{1}{c|}{Краткое содержание документа}\\

\hline

2. Тезаурус\newline инфраструктуры ЕИП РФ&

Тезаурус как предметный словарь, полномерно

охватывающий понятия, определения и~термины ЕИП РФ и~его инфраструктуры.

Нацеленность тезауруса должна способствовать правильной лексической, корпоративной

коммуникации в~инфраструктуре ЕИП РФ, пониманию в~общении и~взаимодействии лиц

и\!/\!или ИР в~рамках ЕИП РФ. Тезаурус в~электронном формате как

действенный инструмент для описания отдельных предметных областей в~рамках ЕИП РФ.

Возможности тезауруса выявлять смысл не только с~по\-мощью определения, но

и~посредством соотнесения слова с~другими понятиями и~их группами в~рамках ЕИП РФ.

Возможности использования тезауруса для наполнения баз знаний систем искусственного

интеллекта в~рамках ЕИП РФ\\

\hline

3. Совокупность стандартов в час\-ти инфраструктуры ЕИП РФ&Правила ведения

нормативной информации трех видов~--- правовой, технической, справочной. Форматы

данных, протоколы обмена информацией, программные интерфейсы, управление

требованиями. Типовые технические решения, стиль программирования, управление

версиями, управление конфигурациями. Порядок испытания и~тестирования систем

Участников~ЕИП РФ. Требования к~документированию со стороны Участников ЕИП РФ\\

\hline

4. Описание про-\newline цес\-сов Участника в~со\-ста\-ве ЕИП РФ, соглашение

с~Участником об\newline

уровне предо\-став\-ле\-ния услуг и регламенты&ИТ-услуги каждого Участника ЕИП РФ и

правила их предо\-став\-ле\-ния другим Участникам в составе ЕИП РФ, регламенты

предоставления и получения ИТ-услуг. Соглашение об уровне предоставления услуги

(\textit{англ}.\

Service Level Agreement~--- SLA), обозначающее формальный договор между заказчиком

услуги и~ее поставщиком, содержащее описание услуги, права и~обязанности сторон

и~согласованный уровень качества предоставления данной услуги\\

\hline

5. Схема информатизации Участника ЕИП РФ &Для каждого Участника ЕИП РФ: кто,

чем и~для чего пользуется Участник, доступность сервисов и~систем в~составных частях

Участника, автоматизированные и~неавтоматизированные участки работы\\

\hline

6. Схема информационных потоков&Точки поступления информационных потоков

(информации) от Участников в~инфраструктуру ЕИП РФ, хранение информационных

потоков (информации), обмен информацией между Участниками, дублирование информации

и~операций по ее вводу\\

\hline

7. Схема взаим\-ной зависимости ИР Участников ЕИП РФ и~их

сервисов &Описание степени зависимости ИР Участников, а~также

их сервисов и~механизмов (критическая, некритическая, слабая, сильная и~др.)\\

\hline

\end{tabular}

\end{center}

\end{table}



     \section{Заключение}



     К настоящему времени, несмотря на целый ряд принятых

     норма\-тив\-но-пра\-во\-вых документов федерального уровня, одной из ключевых проблем, не

позволяющих получить реальные результаты по созданию

ЕИП, является неполнота и несовершенство

     нор\-ма\-тив\-но-тех\-ни\-че\-ско\-го обеспечения процессов разработки,

внедрения и поддержания функционирования ЕИП РФ.



     В статье представлены результаты, направленные на повышение

эффективности процессов формирования ЕИП РФ:

     \begin{enumerate}[(1)]

     \item обоснованы основные принципиальные положения по созданию и

использованию ИР в~составе ЕИП РФ. При этом ЕИП

РФ рассматривается как инфраструктурное образование~--- мегасистема;

     \item обоснованы предложения по перечню

     нор\-ма\-тив\-но-тех\-ни\-че\-ских документов для обеспечения создания

ЕИП РФ как мегасистемы;

     \item сформулировано содержание каждого из предлагаемых документов.

     \end{enumerate}



{\small\frenchspacing

{%\baselineskip=10.8pt

\addcontentsline{toc}{section}{Литература}

\begin{thebibliography}{99}

\bibitem{1-zk}

Концепция формирования и развития единого информационного пространства России и

соответствующих государственных информационных ресурсов. Одобрена решением

Президента РФ 23~ноября 1995~г. №\,Пр-1694.

\bibitem{2-zk}

Электронная Россия: Федеральная целевая программа. Последние изменения 2008~года.

\bibitem{3-zk}

Стратегия развития информационного общества в Российской Федерации. Утверж\-де\-на

Президентом РФ 7~февраля 2008~г. №\,Пр-212.



\bibitem{5-zk} %4

Концепция долгосрочного социально-экономического развития Российской Федерации на

период до 2020~года. Утверждена распоряжением Правительства Российской Федерации от

17~ноября 2008~г. №\,1662-р.



\bibitem{7-zk}  %5

Об организации предоставления государственных и муниципальных услуг: Федеральный

Закон РФ от 27~июля 2010~г. №\,210-ФЗ.



\bibitem{4-zk} %6

Информационное общество (2011--2020~годы): Государственная программа Российской

Федерации. Утверждена распоряжением Правительства РФ №\,1815-р от 20~октября 2010~г.



\bibitem{6-zk} %7

Развитие отрасли информационных технологий: План мероприятий (<<дорожная карта>>).

Утвержден распоряжением Правительства Российской Федерации от 30~декабря 2013~г.

№\,2602-р.



\bibitem{13-zk} %8

Концепция формирования и развития единого информационного пространства России и

соответствующих государственных информационных ресурсов.~--- М.: Информрегистр,

1996.



\bibitem{17-zk} %9

\Au{Данилин А.\,В., Слюсаренко А.\,И.} Архитектура и стратегия. <<Янь>> и~<<Инь>>

информационных технологий предприятия.~--- М.: Ин\-тер\-нет-уни\-вер\-си\-тет

информационных технологий, 2005. 504~с.



\bibitem{10-zk} %10

Об утверждении Типовой программы развития и использования информационных и

телекоммуникационных технологий субъекта Российской Федерации: Распоряжение

Правительства РФ от 3~июля 2007~г. №\,871-р.



\bibitem{9-zk} %11

Об одобрении Концепции формирования в Российской Федерации электронного

правительства до 2010~года: Распоряжение Правительства РФ от 6~мая 2008~г. №\,632-р.



\bibitem{18-zk} %12

\Au{Когаловский М.\,Р., Хохлов Ю.\,Е.} Стандарты XML для электронного правительства.~---

М.: Институт развития информационного общества, 2008. 416~с.



\bibitem{12-zk} %13

Концепция единой системы информационно-справочной поддержки граждан и~организаций

по вопросам взаимодействия с~органами исполнительной влас\-ти и~органами местного

самоуправления с~использованием ин\-фор\-ма\-ци\-он\-но-те\-ле\-ком\-му\-ни\-ка\-ци\-он\-ной сети Интернет.

Одобрена постановлением Правительства РФ от 15~июня 2009~г. №\,478.



\bibitem{14-zk} %14

Правила размещения в федеральных государственных информационных системах

<<Сводный реестр государственных и муниципальных услуг (функций)>> и <<Единый

портал государственных и муниципальных услуг (функций)>> сведений о государственных

и муниципальных услугах (функциях). Утверждены постановлением Правительства РФ от

15~июня 2009~г. №\,478.



\bibitem{11-zk} %15

 Об утверждении Плана перехода на предоставление государственных услуг и исполнение

государственных функций в электронном виде федеральными органами исполнительной

власти: Распоряжение Председателя Правительства РФ от 17~октября 2009~г. №\,1555-р.



\bibitem{15-zk} %16

Сводный перечень первоочередных государственных и муниципальных услуг,

предо\-став\-ля\-емых органами исполнительной власти субъектов Российской Федерации

и~органами местного самоуправления в~электронном виде, а~также услуг, предоставляемых в

электронном виде учреждениям субъектов Российской Федерации и~муниципальными

учреждениями. Утвержден распоряжением Правительства РФ от 17~декабря 2009~года

№\,1993-р. Приложение~1.



\bibitem{8-zk} %17

О~единой вертикально интегрированной государственной автоматизированной

информационной системе <<УПРАВЛЕНИЕ>>: Постановление Правительства РФ от

25~декабря 2009~г. №\,1088.



\bibitem{19-zk} %18

\Au{Штрик А.\,А.} Использование информационно-коммуникационных технологий для

экономического развития и государственного управления в странах современного мира~/\!\!/

Информационные технологии. Приложение, 2009. №\,6. 32~с.

\bibitem{20-zk} %19

\Au{Штрик А.\,А.} Электронные технологии в~деятельности органов государственной власти

России: анализ и~перспективы развития~/\!\!/ Информационные технологии. Приложение,

2009. №\,10. 32~с.

\bibitem{21-zk} %20

\Au{Штрик А.\,А.} Состояние и перспективы формирования информационного общества

в~России до 2015~года~/\!\!/ Информационные технологии. Приложение, 2010. №\,6. 32~с.



\bibitem{16-zk} %21

\Au{Колин К.\,К.} Вызовы XXI века и~стратегические приоритеты развития России~/\!\!/

Стратегические приоритеты, 2014. №\,2. С.~25--42.



\end{thebibliography}

} }





\vspace*{-3pt}



\hfill{\small\textit{Поступила в редакцию 12.09.14}}



\newpage





%\vspace*{9pt}



%\hrule



%\vspace*{2pt}



%\hrule



%\vspace*{9pt}





\def\tit{SOME APPROACHES TO FORMING REGULATORY

AND~TECHNICAL BASE FOR~CREATION OF~UNIFIED INFORMATION SPACE OF~RUSSIA}



\def\titkol{Some approaches to forming regulatory and technical

base for %creation of

unified information space} % of Russia}



\def\aut{A.\,A.~Zatsarinny and E.\,V.~Kiselev}

\def\autkol{A.\,A.~Zatsarinny and E.\,V.~Kiselev}





\titel{\tit}{\aut}{\autkol}{\titkol}



\vspace*{-12pt}



\noindent

Institute of Informatics Problems, Russian Academy of Sciences,

44-2 Vavilov Str., Moscow 119333, Russian

Federation







\def\leftfootline{\small{\textbf{\thepage}

\hfill Sistemy i Sredstva Informatiki~---

Systems and Means of Informatics 2014 vol~24 no\,4}

}%

 \def\rightfootline{\small{Sistemy i Sredstva Informatiki~---

 Systems and Means of Informatics   2014   vol~24   no\,4

\hfill \textbf{\thepage}}}



\Abste{Targets of forming and developing the Russian Federation unified

 information space are determined relying on federal level regulatory and legal

 documents. The article defines the generalized composition of the unified

 information space based on aggregation of information resources and information

 infrastructure of the country. Authors created the definition of a Member of the

 space as a principle producer and customer of information resources.

 The article identifies three main problems of creating and using a unified

 information space considering the real conditions of informatization of the

  country. Based on the systemic approach,

  the article determines general principal theses regarding creation and usage

  of a unified information space. The article formulates the approach and proposals

  for documentation of a unified information space as a~megasystem including seven

  absolutely new types of documents. The results can serve as the basis for further

  investigations concerning creation and interaction of information resources

  associated with the unified information space of the Russian Federation.}



\KWE{unified information space of Russian Federation; information resources;

 information

infrastructure; centralization; integration; information security;

providing telecommunication; systemic approach; megasystem;

regulatory and technical document}







\DOI{10.14357\!/\!08696527140413}



%\vspace*{-24pt}



%\Ack

%\noindent





\vspace*{-2pt}







\renewcommand{\bibname}{\protect\rmfamily References}

%\renewcommand{\bibname}{\large\protect\rm References}



{\small\frenchspacing

{%\baselineskip=10.8pt

\addcontentsline{toc}{section}{References}

\begin{thebibliography}{99}

\bibitem{1-zk-1} %1

PR-1694. Kontseptsiya formirovaniya i razvitiya edinogo informatsionnogo prostranstva Rossii i

sootvetstvuyushchikh gosudarstvennykh informatsionnykh resursov. Odobrena resheniem

Prezidenta Rossiyskoy Federatsii 23~noyabrya 1995~g.\ [A~conception of forming and development of

unified information space of Russia and relevant state information resources. Approved by the

decision of President of the Russian Federation on November~23, 1995].

\bibitem{2-zk-1} %2

Postanovlenie Pravitel'stva No.\,65. 2002. Federal'naya tselevaya programma ``Elektronnaya

Rossiya.'' Poslednie izmeneniya 2008~goda [Resolution of the Government of Russian Federation

No.\,65 of 2009. Federal objective program ``Electronic Russia.'' Including changes of 2008].

\bibitem{3-zk-1} %3

PR-212. Strategiya razvitiya informatsionnogo obshchestva v~Rossiyskoy Federatsii. Utverzhdena

Prezidentom Rossiyskoy Federatsii 7~fevralya 2008~g.\ [Strategy of information society development

in Russian Federation. Approved by President of Russian Federation on February 7, 2008].

\bibitem{5-zk-1} %4

Kontseptsiya dolgosrochnogo sotsial'no-ekonomicheskogo razvitiya Rossiyskoy Federatsii na

period do 2020~goda. Utverzhdena rasporyazheniem Pravitel'stva Rossiyskoy Federatsii

No.\,1662-R ot 17~noyabrya 2008~g.\ [A~conception of long-time social and economic development

of Russian Federation until 2020. Approved by the Decree of the Government of Russian

Federation No.\,1662-R on November 17, 2008].



\bibitem{7-zk-1} %5

Fegeral'nyy Zakon Rossiyskoy Federatsii ot 27.07.2010 No.\,210-FZ

[210-FZ Federal Law of RF dated

27.07.2010]. Ob organizatsii predostavleniya gosudarstvennykh i munitsipal'nykh uslug [Regarding

organization of providing state and municipal services].



\bibitem{4-zk-1} %6

Gosudarstvennaya programma Rossiyskoy Federatsii ``Informatsionnoe obshchestvo (2011-2020

gody).'' Utverzhdena rasporyazheniem Pravitel'stva RF No.\,1815-R ot

20~oktyabrya 2010~g.\ [State

program of Russian Federation ``Information society (2011--2020).'' Approved by the Decree

of the Government of Russian Federation No.\,1815-R on October~20, 2010].





\bibitem{6-zk-1} %7

Razvitie otrasli informatsionnykh tehnologiy:

Plan meropriyatiy (``dorozhnaya karta'').

Utverzhden rasporyazheniem Pravitel'stva Rossiyskoy Federatsii No.\,2602-R ot

30~dekabrya 2013~g.\

[Roadmap ``Information technology industry development.'' Approved by

the Decree of the

Government of Russian Federation No.\,2602-R on December~30, 2013].



\bibitem{13-zk-1} %8

Kontseptsiya formirovaniya i razvitiya edinogo informatsionnogo prostranstva Rossii

i~sootvetstvuyushchikh gosudarstvennykh informatsionnykh resursov [A~conception of forming and

development of unified information space of Russia and relevant state information resources]. 1996.

Moscow: Informregistr Publ.



\bibitem{17-zk-1} %9

\Aue{Danilin, A.\,V., and A.\,I.~Slyusarenko}. 2005. \textit{Arkhitektura i strategiya. ``Yan'\,'' i

``In'\,'' informatsionnykh tekhnologiy predpriyatiya} [Architecture and strategy. Yin and yang of

enterprise information techology]. Moscow: Internet-universitet informatsionnykh tekhnologiy

[Internet-University of information technology]. 504~p.





\bibitem{10-zk-1} %10

Rasporyazhenie pravitelstva RF No.\,871-R ot 3~iyulya 2007~g. Ob utverzhdenii Tipovoy programmy

razvitiya i ispol'zovaniya informatsionnykh i~telekommunikatsionnykh tehnologiy sub''ekta

Rossiyskoy Federatsii [Decree of the Government of Russian Federation No.\,871-R dated

3.07.2007. Regarding approval of typical program of development and using of information and

telecommunication technologies of a Russian Federation subject].



\bibitem{9-zk-1} %11

Rasporyazhenie Pravitelstva RF No.\,632-R ot 6~maya 2008~g. Ob odobrenii Kontseptsii

formirovaniya v~Rossiyskoy Federatsii elektronnogo pravitel'stva do 2010~goda [Decree of the

Government of Russian Federation No.\,632-R dated 6.05.2008. About approval of the Conception

of forming electronic government until 2010].



\bibitem{18-zk-1} %12

\Aue{Kogalovskiy, M.\,R., and Y.\,E.~Hohlov}. 2008. \textit{Standarty XML dlya elektronnogo

pravitel'stva} [XML standards for electronic government]. Moscow: Institut razvitiya

informatsionnogo obshchestva [Institute of information society development]. 416~p.



\bibitem{12-zk-1} %13

Kontseptsiya edinoy sistemy informatsionno-spravochnoy podderzhki grazhdan i organizatsiy po

voprosam vzaimodeystviya s organami ispolnitel'noy vlasti i organami mestnogo samoupravleniya s

ispol'zovaniem informatsionno-telekommmunikatsionnoy seti Internet. Odobrena postanovleniem

Pravitel'stva RF No.\,478 ot 15~iyunya 2009~g.\ [A~conception of unified system of

information-referential support of citizens and organizations in questions of interactions with execution

authorities and local governments using information-telecommunication network Internet.

Approved by the Resolution of the Government of Russian Federation No.\,478 dated 15.06.2009].





\bibitem{14-zk-1} %14

Pravila razmeshcheniya v federal'nykh gosudarstvennykh informatsionnykh sistemakh ``Svodnyy

reestr gosudarstvennykh i munitsipal'nykh uslug (funktsiy)'' i~``Edinyy portal gosudarstvennykh i

munitsipal'nykh uslug (funktsiy)'' svedeniy o~gosudarstvennykh i~munitsipal'nykh

uslugakh

(funktsiyakh). Utverzhdeny postanovleniem Pravitel'stva RF No.\,478 ot

15~iyunya 2009~g.\ [The Rules

of placement in Federal state information systems ``Consolidated register of state and municipal

services (functions)'' and ``Unified portal of state and municipal services (functions)'' of state and

municipal services (functions) data. Approved by the Resolution of the Government of Russian

Federation No.\,478 dated 15.06.2009].



\bibitem{11-zk-1} %15

Rasporyazhenie Predsedatelya Pravitelstva RF No.\,1555-R ot 17~oktyabrya 2009~g. Ob utverzhdenii

Plana perekhoda na predostavlenie gosudarstvennykh uslug i ispolnenie gosudarstvennykh funktsiy v

elektronnom vide federal'nymi organami ispolnitel'noy vlasti [Decree of the Government of Russian

Federation No.\,1555-R dated 17.10.2009. Regarding approval of transition plan to providing of

state services and execution of state functions electronically by federal execution authorities].



\bibitem{15-zk-1} %16

Svodnyy perechen' pervoocherednykh gosudarstvennykh i munitsipal'nykh,

predostav\-lyaemykh organami ispolnitel'noy vlasti sub''ektov Rossiyskoy

Federatsii i~\mbox{organami} mestnogo

samoupravleniya v~elektronnom vide, a takzhe uslug, pre\-dostav\-lya\-emykh v~elektronnom vide

uchrezhdeniyam sub''ektov Rossiyskoy Federatsii i munitsipal'nymi uchrezhdeniyami.

Prilozhenie~1. Utverzhden rasporyazheniem Pravitel'stva RF No.\,1993-R to

17~dekabrya 2009~g.

[Consolidated list of primary state and municipal services provided by execution authorities of

Russian Federation subjects and local governments electronically and services provided to

institutions of Russian Federation subjects and by municipal institutions electronically. Appendix 1.

Approved by the Decree of the Government of Russian Federation No.\,1993-R dated 17.12.2009.]



\bibitem{8-zk-1} %17

Postanovlenie Pravitelstva RF No.\,1088 ot 25~dekabrya 2009~g. O edinoy vertikal'no

integrirovannoy gosudarstvennoy avtomatizirovannoy informatsionnoy sisteme UPRAVLENIE

[Resolution of the Government of Russian Federation No.\,1088 dated 25.12.2009. About unified

vertical integrated state automated information system ``Management''].





\bibitem{19-zk-1} %18

\Aue{Shtrik, A.\,A.} 2009. Ispol'zovanie informatsionno-kommunikatsionnykh tehnologiy dlya

ekonomicheskogo razvitiya i gosudarstvennogo upravleniya v stranakh sovremennogo mira [Using

of information and telecommunication technology for economic development and governance in

modern world countries]. \textit{Prilozhenie k~zhurnalu ``Informatsionnye Tekhnologii''}

[Information Technologies. Appendix] 6. 32~p.

\bibitem{20-zk-1} %19

\Aue{Shtrik, A.\,A.} 2009. Elektronnye tehnologii v deyatel'nosti organov gosudarstvennoy vlasti

Rossii: Analiz i perspektivy razvitiya [Electronic technologies in Russia Government authorities

activities: Analysis and prospect of development]. \textit{Prilozhenie k~zhurnalu ``Informatsionnye

Tehnologii''} [Information Technologies. Appendix] 10. 32~p.

\bibitem{21-zk-1} %20

\Aue{Shtrik, A.\,A.} 2010. Sostoyanie i perspektivy formirovaniya informationnogo

ob\-shche\-st\-va

v~Rossii do 2015 goda [Conditions and prospects of forming information society in Russia until

2015]. \textit{Prilozhenie k zhurnalu ``Informatsionnye Tekhnologii''} [Information Technologies.

Appendix]  6. 32~p.



\bibitem{16-zk-1} %21

\Aue{Kolin, K.\,K.} 2014. Vyzovy XXI~veka i~strategicheskie prioritety razvitiya Rossii  [21st

century demands and strategy priority of Russia development]. \textit{Strategicheskie Prioritety}

[Strategy Priorities] 2:25--42.



\end{thebibliography}

} }





\vspace*{-3pt}



\hfill{\small\textit{Received September 12, 2014}}





\Contr



\noindent

\textbf{Zatsarinnyy Alexander A.} (b.\ 1951)~--- Doctor of Science in technology,

professor, Deputy Director, Institute of Informatics Problems, Russian Academy of

Sciences, 44-2 Vavilov Str., Moscow 119333, Russian Federation; azatsarinny@ipiran.ru



\vspace*{3pt}



\noindent

\textbf{Kiselev Eduard V.} (b.\ 1938)~--- senior scientist, Institute of Informatics Problems,

Russian Academy of Sciences, 44-2 Vavilov Str., Moscow 119333, Russian Federation;

ekiselev@ipiran.ru











 \label{end\stat}



\renewcommand{\bibname}{\protect\rm Литература}









